\documentclass[10pt,a4paper]{amsart}
\usepackage[margin=19mm]{geometry}
\usepackage{amsmath,amssymb,mathtools}
\usepackage[scr=rsfs,cal=euler]{mathalfa}
\usepackage{xfrac}
\usepackage[unicode,hidelinks,bookmarks=false]{hyperref}
\newcommand{\ie}{i.e.\ }
\newcommand{\cf}{cf.\ }
\newcommand{\resp}{respectively\ }
\newcommand{\Subsection}{\subsection}
\providecommand{\nobreakdash}{\nobreak\hbox{-}}
\setlength{\emergencystretch}{2em}
\multlinegap0pt
\usepackage{bm}
\usepackage{bbm}
\usepackage{dsfont}
\usepackage{xspace}
\usepackage{cancel}
\usepackage{mathtools}
\usepackage{amsthm}
\makeatletter
\@ifundefined{theo}{\newtheorem{theo}{Theo}}{}
\@ifundefined{lemm}{\newtheorem{lemm}[theo]{Lemma}}{}
\makeatother
\title[Unboundedness of zero-cycles on higher dimensional Fano mani]{Unboundedness of zero-cycles on higher dimensional Fano manifolds}
\author{Claire Voisin}
\date{Source: arXiv:2511.23080v2 (CC BY 4.0)}
\begin{document}
\maketitle

\noindent\emph{Source and licence.} Unboundedness of zero-cycles on higher dimensional Fano manifolds. Version arXiv:2511.23080v2 (CC BY 4.0), distributed under the Creative Commons Attribution 4.0 International licence, as recorded in the arXiv licence metadata for this exact version and shown on the abstract page https://arxiv.org/abs/2511.23080v2. Full rights notice: licenses/arxiv-2511.23080-CC-BY-4.0.txt.\par\smallskip

\noindent\textit{Editorial scope.} Counted unit: the Lemm whose source label is \texttt{lerankprime} in the article (source0.tex lines 419--428, source label \texttt{lerankprime}), delivered together with the article's own complete original proof of it (source0.tex lines 430--445, one \texttt{proof} environment of 16 lines). No mathematics is written, repaired or re-derived: statement and proof are the article's own text line for line. Required context retained uncounted: none. Every definition, display and label of those lines is retained unchanged. Omitted from this package and not redistributed: 1--418, 429--429, 446--1022. Nothing inside a retained excerpt is omitted or altered: every retained body is delivered exactly as the article's lines are written, so no omission receipt of any kind accompanies this package. Citations: the retained excerpts carry no citation command at all, so no bibliography entry of the article is transcribed and no reference is redistributed. Reference adaptation: none was needed and none was made; every reference inside the delivered unit resolves inside it, so no unresolved reference or citation is produced. Printed slips kept literal: no printed word, symbol, hypothesis or number of the counted statement or of its proof was altered. Layout and class: the article's own class is \texttt{article} with packages including color, amssymb, amsthm, array, amscd, amsfonts, latexsym, url, amsmath, xy; the wrapper is the lane's pinned \texttt{amsart} editorial wrapper at 10pt on A4, which restates the theorem-like environments the retained text needs and adds the article's own macro definitions that the retained text needs (none), with extra wrapper packages (bm, bbm, dsfont, xspace, cancel, mathtools). The delivered numbering is the unit's own. Assets: no retained span contains a figure environment, a graphics-inclusion command, a PDF-page inclusion, a table or a TikZ picture; no image file is copied into this package, the package declares no asset dependency and no image file is redistributed with it. Licence: version arXiv:2511.23080v2 of this article is distributed under the Creative Commons Attribution 4.0 International licence, as recorded in the arXiv licence metadata for this exact version and shown on the abstract page https://arxiv.org/abs/2511.23080v2. A full rights notice is delivered at licenses/arxiv-2511.23080-CC-BY-4.0.txt. Characters: every retained excerpt is pure ASCII, so the delivered engine, XeLaTeX with its default Latin Modern text font, reports no missing character.
\begin{lemm}\label{lerankprime} (i) Let $\phi: V\rightarrow W$ be a linear map, and let $\omega\in\bigwedge^l W^*$. Then ${\rm trk}(\phi^*\omega)\leq {\rm trk}(\omega)$.

(ii) Let $V$ be a vector space of the form $V=V_1\bigoplus \ldots\bigoplus V_N$ and let $\omega_i\in\bigwedge^lV_i^*$. Assume that ${\rm dim}\,V_i=l$ or ${\rm dim}\,V_i=l+1$. Then if $l\geq 2$, the tensor rank of \begin{eqnarray}\label{eqpouromegaprime} \omega=\sum_i{\rm pr}_i^*\omega_i\end{eqnarray} is equal to $N$, assuming $\omega_i\not=0$ for all $i$.

(iii)  In the situation of (ii), assume that either $l=2$ or ${\rm dim}\,V_i=l+2$, so that $\omega_i$ has an even rank $2r_i$ (hence tensor rank $r_i$). Then if $l\geq 2$  and  $ \omega=\sum_i{\rm pr}_i^*\omega_i$,   ${\rm trk}(\omega)=\sum_ir_i$.

(iv) The tensor rank of a $l$-form $\alpha\in \bigwedge^lV^* $ on  a vector space $V$ of dimension $n$ is lower-semicontinuous on $ \bigwedge^lV^* $ equipped with the Zariski topology, assuming $l\leq2$ or $l\geq n-2$.

(v) The tensor rank of a $l$-form is subadditive : ${\rm trk}(\sum_i\omega_i)\leq \sum_i{\rm trk}(\omega_i)$.
\end{lemm}

\begin{proof} (i) follows immediately from the definition of the tensor rank since the pull-back of a decomposable form is decomposable.

For (ii), as any $l$-form on a vector space of dimension $l$ or $l+1$ is decomposable, we get by definition of the tensor rank that a form $\omega$ as in (\ref{eqpouromegaprime}) has tensor rank $\leq N$.  Suppose it has tensor rank $<N$. Then it can be written as
$$\omega=\omega'_1+\ldots+\omega'_{N'}$$
with $N'<N$ and $\omega'_i$ decomposable. It follows that there is a vector subspace $V'_\omega\subset V$ of dimension $\geq {\rm dim}\,V-N'l$ contained in the kernel $V_\omega$ of $\omega$ consisting of vectors $v\in V$ such that $v\lrcorner\omega=0$, namely the space  $V'_\omega$  annihilated by all $\xi_{sr}\in V^*$, where
$\omega'_{s}=\xi_{s1}\wedge\ldots\wedge \xi_{sl}$.
However, for $\omega$ given by (\ref{eqpouromegaprime}), a vector $$v=v_1+\ldots+ v_N\in V=V_1\bigoplus \ldots\bigoplus V_N$$ belongs to the kernel  $V_\omega$ of $\omega$ if and only if \begin{eqnarray}\label{eqvanpourle1}v_i\lrcorner\omega_i=0\,\,{\rm  for\,\, all}\,\, i. \end{eqnarray}
(Note that, as before, this is at this point that we use the condition $l\geq 2$, which implies that each $v_i\lrcorner\omega_i$ has degree $l-1>0$.)
As all $\omega_i$ are nonzero, condition (\ref{eqvanpourle1}) defines a space of codimension $Nl$, which provides  a contradiction.

(iii) is proved similarly as (ii), except that the computation of the tensor rank is different. We use here the fact that when ${\rm dim}\,V_i=l+2$, $\bigwedge^lV_i^*\cong \bigwedge^2V_i$ so the tensor rank of $\omega_i$ is computed as half the rank of $\omega_i$ seen as a $2$-form on $V_i^*$.

(iv) When $n=l$ or $n=l+1$, the tensor rank of any $\omega\in \bigwedge^lV^*$ is $1$ if $\omega\not=0$ and $0$ otherwise. When $n=l+2$, it was already explained  that the tensor rank is half the rank of the corresponding element of $\bigwedge^2V$, and similarly when $l\leq 2$, hence in all these cases, the rank  is lower-semicontinuous.

(v) trivially follows from the definition of the tensor rank.
\end{proof}

\end{document}
